

\section{Background}
\label{sec:pre}

Shilling attacks can achieve both push attacks (promote the target item) and nuke attacks (demote the target item).
Since attackers can easily reverse the goal setting to conduct each attack~\citep{abs-2206-11433}, we consider push attacks in the sequel for simplicity.
In this paper, the \emph{source data} and the \emph{target data} refer to the RS data used for training the attack model and the data in the victim RS, respectively.


\subsection{Related Work}
\label{sec:related}

Early works of shilling attacks rely on heuristics~\citep{GunesKBP14}. 
Recent works~\cite{DeldjooNM21} mostly adopt the idea of adversarial attack~\citep{YuanHZL19}, and they can be categorized into four groups.


\textbf{Optimization methods} study how to model shilling attacks as an optimization task and then use optimization strategies to solve it.
\citet{LiWSV16} assumes the victim RS adopts matrix factorization (MF) and they propose methods PGA and SGLD that directly add the attack goal into the objective of MF.
RevAdv proposed by~\citet{TangWW20} and RAPU proposed by~\citet{ZhangTLSYZG21} model shilling attacks as a bi-level optimization problem.

\textbf{GAN-based methods} adopt Generative Adversarial Network (GAN)~\citep{GoodfellowPMXWOCB14} to construct fake user profiles. The generator models the data distribution of real users and generates real-like data, while the discriminator is responsible for identifying the generated fake users. 
Along this direction, a large number of methods have sprung up: TrialAttack~\citep{WuLGZC21}, Leg-UP~\citep{abs-2206-11433}, DCGAN~\citep{Christakopoulou19}, AUSH~\citep{LinC0XLY20}, RecUP~\citep{ZhangCZWL21}, to name a few. 

\textbf{RL-based methods} query the RS to get feedback on the attack. Then, Reinforcement Learning (RL)~\citep{KaelblingLM96} is used to adjust the attack. Representative works include PoisonRec~\citep{SongLHWLLG20}, LOKI~\citep{ZhangLD020} and CopyAttack~\citep{FanDZ0LWT021}. 

\textbf{KD-based methods} leverage Knowledge Distillation (KD)~\citep{GouYMT21} to narrow the gap between the surrogate RS and the victim RS. The surrogate RS is used to mimic the victim RS when the prior knowledge is not available. 
Model Extraction Attack proposed by \citet{YueHZM21} falls in this category.


\subsection{Analysis of Existing Works}
\label{sec:analysis}

We review existing shilling attack approaches and summarize their \emph{characteristics} in Tab.~\ref{tab:pre}:
\begin{itemize}
	\item \textbf{Data Knowledge}: Some methods assume the complete/partial target data is exposed to attackers. A practical attack method should use as less target data as possible.
	\item \textbf{RS Parameter Knowledge}: Some methods require the knowledge of the learned parameters of the victim RS. Such information is typically not available. 
	\item \textbf{RS Architecture Knowledge}: Some methods require the knowledge of the architecture of the victim RS. Such information is typically not available. 
	\item \textbf{Train with a surrogate RS}: Use a surrogate RS to train the attacker to avoid the prior knowledge of the victim RS.
	\item \textbf{Require multiple queries}: Query the victim RS multiple times and adjust fake  profiles according to the feedback.
	\item \textbf{Cross-domain Attack}: Use the information in one RS domain to attack another RS domain, e.g., train on the book data in Amazon RS and then attack video items in Amazon RS. Source and target domains share users and/or items.
	\item \textbf{Cross-system Attack}: Use the information in one RS to attack another RS, e.g., train on the Yelp RS and then attack Amazon RS. Source RS and target RS may not share users and/or items.
\end{itemize}

\begin{figure*}[t]
\centering
\includegraphics[width=1\linewidth]{framework.pdf}
\caption{Overview of \ours.}
\label{fig:framework}
\end{figure*}

From Tab.~\ref{tab:pre}, we can see that 
none of the existing methods have all the three properties illustrated in Sec.~\ref{sec:intro}.
In other words, \emph{there is still no real practical shilling attack method}.
Particularly, we do not find any method that has Property 2 and can achieve cross-system attack. 
Property 2 partially manifests in CopyAttack.
But CopyAttack assumes the source data and the target data share items and it is only able to achieve cross-domain attack.
Model Extraction Attack considers a data-free setting and uses limited queries ($c$ queries) to close the gap between the surrogate RS and the victim RS.
But its idea only works on sequential RS and the number of required queries is hard to pre-defined.


In summary, based on Tab.~\ref{tab:pre}, we can conclude that our method \ours (illustrated in the next section) has all of the three properties that a practical shilling attack method should have. 
And it is able to achieve cross-system attack, a difficult but practical setting of real shilling attacks.
